% math0 manuscript. Build with: latexmk -pdf paper.tex
\documentclass[11pt]{amsart}
\usepackage[margin=1.1in]{geometry}
\usepackage{amsmath,amssymb,amsthm,mathtools}
\usepackage{url}
\usepackage{hyperref}
\usepackage{xcolor}
\usepackage{listings}
\usepackage{newunicodechar}

\lstdefinelanguage{lean}{
  keywords={theorem,lemma,def,abbrev,example,fun,by,have,show,from,let,in,if,then,else,match,with,
    namespace,end,open,import,section,variable,instance,structure,where,calc,obtain,intro,exact,
    refine,rcases,rintro,simp,omega,decide,norm_num,linarith,nlinarith,positivity,induction,cases,
    private,rw,constructor,right,unfold,by_cases},
  sensitive=true, comment=[l]{--}, morecomment=[s]{/-}{-/}, string=[b]"}
\lstset{language=lean, basicstyle=\ttfamily\footnotesize, columns=fullflexible, keepspaces=true, breaklines=true,
  keywordstyle=\color{blue!60!black}\bfseries, commentstyle=\color{gray}, frame=single,
  extendedchars=true, inputencoding=utf8,
  literate={ℕ}{{$\mathbb{N}$}}1 {ℤ}{{$\mathbb{Z}$}}1 {∀}{{$\forall$}}1 {∃}{{$\exists$}}1
    {↔}{{$\leftrightarrow$}}1 {∧}{{$\wedge$}}1 {∨}{{$\vee$}}1 {→}{{$\to$}}1 {≤}{{$\le$}}1
    {≠}{{$\ne$}}1 {∈}{{$\in$}}1 {·}{{$\cdot$}}1}
% Unicode commonly found in Lean sources.
\newunicodechar{ℕ}{\ensuremath{\mathbb{N}}}
\newunicodechar{ℤ}{\ensuremath{\mathbb{Z}}}
\newunicodechar{ℚ}{\ensuremath{\mathbb{Q}}}
\newunicodechar{ℝ}{\ensuremath{\mathbb{R}}}
\newunicodechar{ℂ}{\ensuremath{\mathbb{C}}}
\newunicodechar{∀}{\ensuremath{\forall}}
\newunicodechar{∃}{\ensuremath{\exists}}
\newunicodechar{→}{\ensuremath{\to}}
\newunicodechar{←}{\ensuremath{\leftarrow}}
\newunicodechar{↔}{\ensuremath{\leftrightarrow}}
\newunicodechar{↦}{\ensuremath{\mapsto}}
\newunicodechar{¬}{\ensuremath{\neg}}
\newunicodechar{∧}{\ensuremath{\wedge}}
\newunicodechar{∨}{\ensuremath{\vee}}
\newunicodechar{≤}{\ensuremath{\le}}
\newunicodechar{≥}{\ensuremath{\ge}}
\newunicodechar{≠}{\ensuremath{\ne}}
\newunicodechar{∈}{\ensuremath{\in}}
\newunicodechar{∉}{\ensuremath{\notin}}
\newunicodechar{⊆}{\ensuremath{\subseteq}}
\newunicodechar{∩}{\ensuremath{\cap}}
\newunicodechar{∪}{\ensuremath{\cup}}
\newunicodechar{∑}{\ensuremath{\sum}}
\newunicodechar{∏}{\ensuremath{\prod}}
\newunicodechar{⟨}{\ensuremath{\langle}}
\newunicodechar{⟩}{\ensuremath{\rangle}}
\newunicodechar{·}{\ensuremath{\cdot}}
\newunicodechar{∘}{\ensuremath{\circ}}
\newunicodechar{⁻}{\ensuremath{^{-}}}
\newunicodechar{¹}{\ensuremath{^{1}}}
\newunicodechar{²}{\ensuremath{^{2}}}
\newunicodechar{₀}{\ensuremath{_{0}}}
\newunicodechar{₁}{\ensuremath{_{1}}}
\newunicodechar{₂}{\ensuremath{_{2}}}
\newunicodechar{α}{\ensuremath{\alpha}}
\newunicodechar{β}{\ensuremath{\beta}}
\newunicodechar{ε}{\ensuremath{\varepsilon}}
\newunicodechar{δ}{\ensuremath{\delta}}
\newunicodechar{π}{\ensuremath{\pi}}
\newunicodechar{σ}{\ensuremath{\sigma}}
\newunicodechar{φ}{\ensuremath{\varphi}}
\newunicodechar{∞}{\ensuremath{\infty}}
\newunicodechar{⌊}{\ensuremath{\lfloor}}
\newunicodechar{⌋}{\ensuremath{\rfloor}}
\newunicodechar{⌈}{\ensuremath{\lceil}}
\newunicodechar{⌉}{\ensuremath{\rceil}}
\newunicodechar{∣}{\ensuremath{\mid}}
\newunicodechar{≡}{\ensuremath{\equiv}}
\newunicodechar{×}{\ensuremath{\times}}
\newunicodechar{⧸}{\ensuremath{/}}
\newunicodechar{‖}{\ensuremath{\|}}
\newunicodechar{√}{\ensuremath{\surd}}
\newunicodechar{∅}{\ensuremath{\emptyset}}
\newunicodechar{⦃}{\ensuremath{\{\!|}}
\newunicodechar{⦄}{\ensuremath{|\!\}}}
\newunicodechar{𝓝}{\ensuremath{\mathcal{N}}}
\newunicodechar{≫}{\ensuremath{\gg}}
\newunicodechar{≪}{\ensuremath{\ll}}

\theoremstyle{plain}
\newtheorem{theorem}{Theorem}
\newtheorem{lemma}[theorem]{Lemma}
\newtheorem{proposition}[theorem]{Proposition}
\newtheorem{corollary}[theorem]{Corollary}
\theoremstyle{definition}
\newtheorem{definition}[theorem]{Definition}
\newtheorem{remark}[theorem]{Remark}

\emergencystretch=3em
% Lean identifiers in running text (breakable at dots and underscores).
\expandafter\def\expandafter\UrlBreaks\expandafter{\UrlBreaks\do\_\do\.}
\DeclareUrlCommand\decl{\urlstyle{tt}}
\newcommand{\N}{\mathbb{N}}
\newcommand{\Z}{\mathbb{Z}}
\newcommand{\rev}{\operatorname{rev}}
\newcommand{\maxdig}{\operatorname{maxdig}}
\newcommand{\val}{\operatorname{val}}

\title[Fixed points of reversal minus largest digit]{The only fixed points of $x\mapsto\rev(x)-\maxdig(x)$ are $0$ and $89$: Earls' conjecture in OEIS A097154}
\author{math0}
\address{Produced by Claude (an AI model developed by Anthropic) in the math0 project}
\date{October 19, 2026}

\begin{document}
\begin{abstract}
For a nonnegative integer $x$ let $\rev(x)$ be the number obtained by reversing the decimal digits of $x$, and let $\maxdig(x)$ be the largest decimal digit of $x$. We prove that the map $f(x)=\rev(x)-\maxdig(x)$ has exactly two fixed points, $0$ and $89$, which was stated as a conjecture by Jason Earls in OEIS A097154 (2004) and repeated in A097155 and A097157.
The proof combines the congruence $\rev(x)\equiv x\pmod 9$, which forces $\maxdig(x)=9$ at a nonzero fixed point, with a comparison of the leading and units digits.
To our knowledge the result is new, although the argument is elementary; the OEIS entries still list the statement as a conjecture.
The theorem is verified in Lean~4 against a statement that math0 formalized from the OEIS entry, reviewed for fidelity and frozen before any proof attempt, whose definitions are checked by the kernel against the listed OEIS terms; this statement is not part of the formal-conjectures repository.
\end{abstract}
\maketitle

\section{Introduction}

In August 2004 Jason Earls contributed to the OEIS a group of sequences A097153--A097158 about the map
\[
  f(x)=\rev(x)-\maxdig(x)\qquad(x\ge0),
\]
where $\rev(x)$ is the decimal reversal of $x$ (sequence A004086 \cite{A004086}, with leading zeros of the reversed string dropped, so that $\rev(120)=21$) and $\maxdig(x)$ is the largest decimal digit of $x$. The values $f(n)$ form A097153 \cite{A097153}. A097154 \cite{A097154} gives the number of iterations of $f$ needed to bring $n$ to $0$, or $-1$ if $0$ is never reached; for example $12\mapsto19\mapsto82\mapsto20\mapsto0$, so that $a(12)=4$. The first value $-1$ is $a(89)$, because $f(89)=98-9=89$. A097155 \cite{A097155} lists the numbers whose orbit reaches $89$, A097157 \cite{A097157} the number of iterations they need to reach $89$, A097156 \cite{A097156} the least number that needs exactly $n$ iterations to reach $0$, and A097158 \cite{A097158} the numbers that set a new record for the number of iterations needed to reach $0$.
The comments of A097154, A097155 and A097157 all contain the sentence ``It is conjectured that 0 and 89 are the only fixed points of f.'' We prove this (Theorem~\ref{thm:main}).

\subsection*{Method}
If $f(x)=x$ then $\rev(x)=x+\maxdig(x)$. Reversal preserves the digit sum, so $\rev(x)\equiv x\pmod 9$, and therefore $9\mid\maxdig(x)$; for $x\ne0$ this gives $\maxdig(x)=9$. It remains to solve $\rev(x)=x+9$ with largest digit $9$. Writing $x$ as units digit $b$, middle block $M$ and leading digit $a$, the equation modulo $10$ leaves two cases for $(a,b)$; one is impossible by size, and the other forces the middle block to be empty, so that $x=10a+(a+1)$, and the condition on the largest digit gives $x=89$.

\subsection*{Status and prior work}
The local copy of the OEIS data used by math0, which the referees cross-checked against the current OEIS data files, has A097154 at revision \#4 (December~15, 2017), A097155 at revision \#11 (August~29, 2025) and A097157 at revision \#3 (December~15, 2017). All three still state the conjecture, and none has a comment, formula, link or edit note recording a proof or a counterexample; the only edit note in A097155 is a correction of its data by Harvey P.~Dale (2011). The related entries A097153, A097156 and A097158, with extensions by Robert G.~Wilson~v and Klaus Brockhaus, are computational, and no other OEIS entry in that copy refers to A097153--A097158.
Web searches by the math0 referees (for the definition of the map, for phrasings such as ``reversal minus largest digit'', and in the SeqFan archives and question-and-answer sites) found no proof; they found only the OEIS pages themselves and work on other digit-reversal maps.
The statement is not among the OEIS statements of the formal-conjectures repository \cite{FC}.
We therefore classify the result as \emph{new}. The argument uses only congruences modulo $9$ and $10$ and a size comparison, and it may well have been noticed without being recorded.

We note that the conjecture is about fixed points in the literal sense, $f(x)=x$; A097155 calls $89$ ``the fixed point 89''. We do not address the stronger question whether $f$ has periodic orbits of length at least $2$, that is, whether every orbit ends at $0$ or at $89$ (Remark~\ref{rem:cycles}).

\section{Statement}\label{sec:statement}

\subsection*{Notation}
For a finite list $\ell=(e_0,\dots,e_{n-1})$ of natural numbers we write $|\ell|=n$,
\[
  \val(\ell)=\sum_{i=0}^{n-1}e_i\,10^i,\qquad \ell^R=(e_{n-1},\dots,e_0),\qquad \mu(\ell)=\max(0,e_0,\dots,e_{n-1}).
\]
In Lean these are \texttt{Nat.ofDigits 10 $\ell$}, \texttt{$\ell$.reverse} and \texttt{$\ell$.foldr max 0}. We use $\val(())=0$, $\val(e::\ell)=e+10\val(\ell)$ and $\val(\ell\mathbin{+\!\!+}\ell')=\val(\ell)+10^{|\ell|}\val(\ell')$, where $e::\ell$ is the list with first entry $e$ followed by $\ell$ and $\mathbin{+\!\!+}$ is concatenation.
For $x\in\N$ let $\delta(x)$ be the list of decimal digits of $x$, least significant first (\texttt{Nat.digits 10 x}). Thus $\delta(0)=()$, and for $x\ge1$ the list $\delta(x)$ is nonempty, its entries lie in $\{0,\dots,9\}$, its last entry (the leading digit) is nonzero, and $\val(\delta(x))=x$.

\begin{definition}\label{def:f}
For $x\in\N$ put
\[
  \rev(x)=\val\bigl(\delta(x)^R\bigr),\qquad \maxdig(x)=\mu\bigl(\delta(x)\bigr),\qquad f(x)=\rev(x)-\maxdig(x).
\]
\end{definition}

So $\rev(x)$ is the number whose decimal digits are those of $x$ in reverse order, with leading zeros dropped, and $\rev(0)=\maxdig(0)=0$. By Lemma~\ref{lem:max}(c) below, $\maxdig(x)\le\rev(x)$, so $f$ maps $\N$ to $\N$ and the subtraction is exact. For example $f(10)=1-1=0$ and $f(89)=98-9=89$.

\begin{theorem}\label{thm:main}
For every $x\in\N$, $f(x)=x$ if and only if $x=0$ or $x=89$.
\end{theorem}

\subsection*{The OEIS conjecture}
The entry A097154 \cite{A097154}, contributed by Jason Earls on August~1, 2004, is named ``Number of iterations of f(x) = reverse(x) - maxdigit(x) needed to reach zero, or -1 if zero is never reached, where maxdig(n) is the largest digit of n'', with offset $1$ and keywords \texttt{base}, \texttt{easy}, \texttt{sign}. Its COMMENTS section reads, verbatim (the comment carries no separate signature):
\begin{quote}
{\small\raggedright\texttt{a(89)=-1 because 89 is the first number that will not reach zero: 98-9=89. It is conjectured that 0 and 89 are the only fixed points of f. See A097155 for numbers that reach 89.}\par}
\end{quote}
The same conjecture appears in the comments of A097155 and A097157, also contributed by Earls on August~1, 2004.

\subsection*{Lean formalization}
The statement was formalized by math0 in the file \path{lean/Math0/Statements/OEIS/A097154.lean}. It was reviewed for fidelity and frozen before any proof attempt; its SHA-256 hash is \texttt{1d7d144b\ldots cbb8}. The definitions of the map and the statement read as follows (verbatim; the import, docstrings, blank lines and all other declarations are omitted):
\begin{lstlisting}
namespace Math0.Statements.A097154
def rev (x : ℕ) : ℕ := Nat.ofDigits 10 (Nat.digits 10 x).reverse
def maxDigit (x : ℕ) : ℕ := (Nat.digits 10 x).foldr max 0
def f (x : ℕ) : ℕ := rev x - maxDigit x
def statement : Prop :=
  ∀ x : ℕ, f x = x ↔ x = 0 ∨ x = 89
end Math0.Statements.A097154
\end{lstlisting}
These are the definitions $\rev$, $\maxdig$ and $f$ of Definition~\ref{def:f}, and \decl{statement} is Theorem~\ref{thm:main}; the biconditional says both that $0$ and $89$ are fixed points and that there are no others. The subtraction in \decl{f} is truncated subtraction on $\N$, but the file proves that it never truncates, and the same file also encodes the sequence A097154 itself. Its other declarations include the following (statements shown, proofs omitted):
\begin{lstlisting}
theorem maxDigit_le_rev (x : ℕ) : maxDigit x ≤ rev x
theorem f_eq_int (x : ℕ) : (f x : ℤ) = (rev x : ℤ) - (maxDigit x : ℤ)
theorem f_examples :
    f 12 = 19 ∧ f 19 = 82 ∧ f 82 = 20 ∧ f 20 = 0 ∧ f 0 = 0 ∧ rev 89 = 98 ∧ maxDigit 89 = 9 ∧
      f 89 = 89
def search : ℕ → ℕ → ℕ → ℤ
  | x, k, 0 => if x = 0 then (k : ℤ) else -1
  | x, k, fuel + 1 => if x = 0 then (k : ℤ) else search (f x) (k + 1) fuel
def a (n : ℕ) : ℤ := search n 0 (10 ^ (Nat.digits 10 n).length)
open Classical in
theorem a_spec (n : ℕ) : a n = if h : ∃ k, f^[k] n = 0 then (Nat.find h : ℤ) else -1
\end{lstlisting}
By \decl{maxDigit_le_rev} and \decl{f_eq_int}, \decl{f} agrees with the integer-valued map of the OEIS entry, so the truncated subtraction creates no junk values.
The sequence takes the value $-1$, so \decl{a} has values in $\Z$. It follows the orbit of $n$ for at most $10^d$ steps, where $d$ is the number of digits of $n$, which makes it computable; \decl{a_spec} shows that this bound loses nothing, so that \decl{a} is exactly the OEIS definition (the least $k$ with $f^k(n)=0$, and $-1$ if there is none). Its proof shows that $f$ maps $\{0,\dots,10^d-1\}$ into itself and that the iterates before the first $0$ are pairwise distinct, and then applies the pigeonhole principle.
Two lemmas proved by \texttt{decide +kernel} compare the definitions with the entry: \decl{a_data} checks all $127$ listed terms $a(1),\dots,a(127)$ of A097154, including $a(89)=-1$, and \decl{f_examples} checks the example $12\mapsto19\mapsto82\mapsto20\mapsto0$ of the entry, together with $f(0)=0$ and $f(89)=98-9=89$.
The reversal convention (trailing zeros of $x$ disappear) is the one the data require, since $a(10)=a(20)=\dots=1$.
The statement quantifies over all $x\in\N$, including $x=0$, which the conjecture names as a fixed point, and adds no hypotheses.

\section{Proof}\label{sec:proof}

The proof follows the Lean file \path{lean/Math0/Frontier/OEIS/A097154.lean}; the name of the corresponding Lean declaration is given with each step. Lemma~\ref{lem:mem} and parts (b) and (c) of Lemma~\ref{lem:max} are proved in the statement file and reused; part (a) of Lemma~\ref{lem:max} is proved in the proof file.

\begin{lemma}[\decl{le_ofDigits_of_mem}]\label{lem:mem}
If $e$ is an entry of the list $\ell$, then $e\le\val(\ell)$.
\end{lemma}

\begin{proof}
By induction on $\ell$. The empty list has no entries. If $\ell=e'::\ell'$, then $\val(\ell)=e'+10\val(\ell')$; either $e=e'\le\val(\ell)$, or $e$ is an entry of $\ell'$ and $e\le\val(\ell')\le\val(\ell)$ by induction.
\end{proof}

\begin{lemma}[\decl{le_foldr_max}, \decl{foldr_max_le}, \decl{maxDigit_le_rev}]\label{lem:max}
Let $\ell$ be a list of natural numbers.
\begin{enumerate}
\item[(a)] Every entry $e$ of $\ell$ satisfies $e\le\mu(\ell)$.
\item[(b)] If $c\in\N$ and every entry of $\ell$ is at most $c$, then $\mu(\ell)\le c$.
\item[(c)] For every $x\in\N$, $\maxdig(x)\le\rev(x)$.
\end{enumerate}
\end{lemma}

\begin{proof}
Parts (a) and (b) follow by induction on $\ell$ from $\mu(())=0$ and $\mu(e::\ell)=\max(e,\mu(\ell))$. For (c), every entry of $\delta(x)$ is an entry of $\delta(x)^R$, hence at most $\val(\delta(x)^R)=\rev(x)$ by Lemma~\ref{lem:mem}; now apply (b) with $c=\rev(x)$.
\end{proof}

\begin{lemma}[\decl{ofDigits_reverse_mod}]\label{lem:mod9}
For every list $\ell$ of natural numbers, $\val(\ell^R)\equiv\val(\ell)\pmod 9$.
\end{lemma}

\begin{proof}
Since $10\equiv1\pmod9$, we have $\val(\ell)=\sum_i e_i10^i\equiv\sum_i e_i\pmod 9$ (in Lean, \decl{Nat.ofDigits_modEq} and \decl{Nat.ofDigits_one}), and likewise for $\ell^R$. The entries of $\ell$ and $\ell^R$ have the same sum (\decl{List.sum_reverse}).
\end{proof}

\begin{lemma}[\decl{ofDigits_eq_zero}]\label{lem:zero}
If every entry of $\ell$ is $0$, then $\val(\ell)=0$.
\end{lemma}

\begin{proof}
Immediate by induction on $\ell$.
\end{proof}

We also use that $\val(\ell)<10^{|\ell|}$ whenever all entries of $\ell$ are at most $9$ (\decl{Nat.ofDigits_lt_base_pow_length} in Mathlib). The heart of the proof is a statement about digit lists.

\begin{proposition}[\decl{key}]\label{prop:key}
Let $\ell$ be a nonempty list with entries in $\{0,\dots,9\}$ whose last entry is nonzero, and suppose that
\begin{equation}\label{eq:hyp}
  \val(\ell^R)=\val(\ell)+\mu(\ell).
\end{equation}
Then $\ell=(9,8)$.
\end{proposition}

\begin{proof}
Put $c=\mu(\ell)$. All entries of $\ell$ are at most $9$, so $c\le9$ by Lemma~\ref{lem:max}(b).

\emph{One entry.} If $\ell=(b)$, then $\ell^R=\ell$, so \eqref{eq:hyp} gives $c=0$. By Lemma~\ref{lem:max}(a), $b\le c=0$, contradicting the assumption that the last entry $b$ is nonzero.

\emph{At least two entries.} Then $\ell=b::(M\mathbin{+\!\!+}(a))$ (\decl{List.eq_nil_or_concat'}), where $b$ is the first entry (the units digit), $a$ is the last entry (the leading digit) and $M$ is a possibly empty list of length $k$. We have $1\le a\le9$, $0\le b\le9$, and $a\le c$, $b\le c$ by Lemma~\ref{lem:max}(a). Put
\[
  m=\val(M),\qquad r=\val(M^R),\qquad P=10^k.
\]
Since $\ell^R=a::(M^R\mathbin{+\!\!+}(b))$ and $|M^R|=k$, the rules for $\val$ give
\begin{equation}\label{eq:vals}
  \val(\ell)=b+10m+10Pa,\qquad \val(\ell^R)=a+10r+10Pb.
\end{equation}
We record three facts.
\begin{enumerate}
\item[(i)] $m<P$ and $r<P$, because the entries of $M$ and of $M^R$ are at most $9$.
\item[(ii)] $P\ge1$, and $P\ge10$ if $k\ge1$.
\item[(iii)] If $r=0$ then $m=0$. Indeed, each entry of $M$ is an entry of $M^R$, hence at most $r=0$ by Lemma~\ref{lem:mem}, and Lemma~\ref{lem:zero} gives $m=0$.
\end{enumerate}
By \eqref{eq:vals}, the hypothesis \eqref{eq:hyp} reads
\begin{equation}\label{eq:main}
  a+10r+10Pb=b+10m+10Pa+c.
\end{equation}

\emph{Step 1: $c=9$.} By Lemma~\ref{lem:mod9}, $c=\val(\ell^R)-\val(\ell)\equiv0\pmod 9$. Since $1\le a\le c\le 9$, we get $c=9$.

\emph{Step 2: $b=a+1$, or $a=9$ and $b=0$.} Reducing \eqref{eq:main} modulo $10$ gives $a\equiv b+9\pmod{10}$. If $b\ge1$ then $a\equiv b-1\pmod{10}$ with $0\le a,\,b-1\le9$, so $a=b-1$, that is, $b=a+1$. If $b=0$ then $a\equiv9\pmod{10}$, so $a=9$.

\emph{Step 3: the case $a=9$, $b=0$ is impossible.} Then \eqref{eq:main} becomes $9+10r=10m+90P+9$, that is $r=m+9P\ge P$, contradicting (i).

\emph{Step 4: the case $b=a+1$ gives $\ell=(9,8)$.} Then \eqref{eq:main} becomes $a+10r+10Pa+10P=a+1+10m+10Pa+9$, that is,
\[
  r+P=m+1.
\]
By (i), $m+1\le P$, so $r=0$. By (iii), $m=0$, so $P=1$, and by (ii), $k=0$; thus $M$ is empty and $\ell=(a+1,a)$. Both entries of $\ell$ are at most $a+1$, so $9=c\le a+1$ by Lemma~\ref{lem:max}(b), while $a+1=b\le9$. Hence $a=8$ and $\ell=(9,8)$.
\end{proof}

In Lean, $m$, $r$ and $P$ are replaced by variables (\texttt{generalize}) after (i)--(iii) and \eqref{eq:vals} have been established, and the arithmetic in Steps~1--4 is done by \texttt{omega}, a decision procedure for linear arithmetic over $\N$ and $\Z$; the products $Pa$ and $Pb$ are treated as atoms, and the identity $P(a+1)=Pa+P$ is supplied by \texttt{ring}. The bound $c\le a+1$ at the end of Step~4 comes from Lemma~\ref{lem:max}(b) (\decl{foldr_max_le}).

\begin{proof}[Proof of Theorem~\ref{thm:main} (\decl{result})]
If $x=0$ or $x=89$, then $f(x)=x$ by direct computation: $f(0)=0-0=0$ and $f(89)=98-9=89$ (these are two conjuncts of \decl{f_examples}, proved by evaluation in the kernel).

Conversely, let $f(x)=x$ and suppose $x\ne0$. By Lemma~\ref{lem:max}(c) the subtraction in $f(x)$ is exact, so $\rev(x)=x+\maxdig(x)$. Apply Proposition~\ref{prop:key} to $\ell=\delta(x)$: it is nonempty because $x\ne0$ (\decl{Nat.digits_ne_nil_iff_ne_zero}), its entries are at most $9$ (\decl{Nat.digits_lt_base}), its last entry is nonzero (\decl{Nat.getLast_digit_ne_zero}), and \eqref{eq:hyp} holds because $\val(\delta(x))=x$ (\decl{Nat.ofDigits_digits}). Hence $\delta(x)=(9,8)$, and $x=\val(\delta(x))=9+10\cdot8=89$.
\end{proof}

\begin{remark}[a by-product; not formalized]\label{rem:plus9}
In Steps~2--4 the hypothesis \eqref{eq:hyp} enters only through \eqref{eq:main} with $c=9$, except in the last sentence of Step~4. Running the same argument for the equation $\rev(x)=x+9$ with $x\ge1$ shows that its solutions are exactly $x=11a+1$ with $1\le a\le8$, that is $12,23,34,45,56,67,78,89$; conversely $\rev(11a+1)=10(a+1)+a=(11a+1)+9$ for these $a$. Step~1 shows that every fixed point $x\ne0$ has $\maxdig(x)=9$, and among these eight numbers only $89$ has largest digit $9$. (The two-digit case of the equation $\rev(x)=x+9$ is a familiar puzzle; the point here is only that no longer solutions exist.) This remark is not part of the Lean development; we checked the list of solutions numerically for $x<10^6$.
\end{remark}

\begin{remark}[periodic orbits; not addressed]\label{rem:cycles}
Since $f$ maps $\{0,\dots,10^d-1\}$ into itself (this is the lemma \decl{f_lt} of the statement file), every orbit of $f$ is eventually periodic. Theorem~\ref{thm:main} determines the periodic orbits of length $1$. We have checked by computer that every orbit starting below $10^6$ ends at $0$ or at $89$, but we do not know whether $f$ has periodic orbits of length at least $2$; this question is not part of the conjecture.
\end{remark}

\section{Machine verification}\label{sec:lean}

\subsection*{Files and declarations}
The proof is in the file \path{lean/Math0/Frontier/OEIS/A097154.lean} of the math0 repository, \url{https://github.com/Foundation-0/math0}. It imports the frozen statement file \path{lean/Math0/Statements/OEIS/A097154.lean} (Section~\ref{sec:statement}) and the audit commands in \path{lean/Math0/Audit.lean}. The main declaration is \decl{Math0.Frontier.A097154.result}, of type \decl{Math0.Statements.A097154.statement}. The helper lemmas \decl{le_foldr_max}, \decl{ofDigits_eq_zero}, \decl{ofDigits_reverse_mod} and \decl{key} (Lemmas~\ref{lem:max}(a), \ref{lem:zero}, \ref{lem:mod9} and Proposition~\ref{prop:key}) are private; the lemmas \decl{le_ofDigits_of_mem}, \decl{foldr_max_le}, \decl{maxDigit_le_rev} and \decl{f_examples} come from the statement file. The core lemma, the main theorem and the final command read as follows (verbatim; proofs omitted):
\begin{lstlisting}[basicstyle=\ttfamily\scriptsize]
private theorem key (L : List ℕ) (hL : ∀ d ∈ L, d < 10) (hlast : ∀ h : L ≠ [], L.getLast h ≠ 0)
    (heq : Nat.ofDigits 10 L.reverse = Nat.ofDigits 10 L + L.foldr max 0) (hne : L ≠ []) :
    L = [9, 8]
theorem result : statement
#math0_settles Math0.Frontier.A097154.result Math0.Statements.A097154.statement
\end{lstlisting}
Here \decl{key} is Proposition~\ref{prop:key}, and the proof of \decl{result} is the proof of Theorem~\ref{thm:main} given above, step by step. The proof file uses no \texttt{native\_decide}, \texttt{sorry} or \texttt{set\_option}; its only use of \texttt{decide} evaluates $\val((9,8))=89$. The lemma \decl{f_examples} in the statement file is proved by \texttt{decide +kernel}, which evaluates in the kernel and adds no axioms.

\subsection*{Audit}
The command \texttt{\#math0\_settles}, defined in \path{lean/Math0/Audit.lean} for statements that math0 formalizes itself, fails unless (1) \decl{statement} is a definition of type \texttt{Prop} without universe parameters; (2) \decl{statement} and \decl{result} use no axioms besides \decl{propext}, \decl{Classical.choice} and \decl{Quot.sound} (this rejects \decl{sorryAx}, \decl{Lean.ofReduceBool} from \texttt{native\_decide}, and custom axioms); and (3) the type of \decl{result} is syntactically the statement or its negation, up to metadata annotations and without unfolding. Here it reports
\begin{lstlisting}[basicstyle=\ttfamily\tiny]
math0 audit: 'Math0.Statements.A097154.statement' depends on axioms [Classical.choice, Quot.sound, propext]
math0 audit: 'Math0.Frontier.A097154.result' depends on axioms [Classical.choice, Quot.sound, propext]
math0 statement check: 'Math0.Frontier.A097154.result' proves 'Math0.Statements.A097154.statement'
\end{lstlisting}

\subsection*{Toolchain and reproduction}
We use Lean~4 \cite{Lean4} \texttt{v4.33.1} and Mathlib \cite{Mathlib} \texttt{v4.33.1}, pinned in \path{lean/}; the repository also pins formal-conjectures \cite{FC} at commit \texttt{c77fdd4}, which this result does not use.
To reproduce, run \path{scripts/verify.sh} from the repository root: it fetches the Mathlib build cache, builds all math0 results with \texttt{lake build} (re-running every audit command), checks that the pinned \texttt{formal\_conjectures} and \texttt{mathlib} checkouts are unmodified, and reports the number of passed statement checks. The single file can be checked with \texttt{lake env lean Math0/Frontier/OEIS/A097154.lean} in the directory \path{lean}; this takes about twenty seconds on our machine once Mathlib and the statement file have been built.

{\small
\begin{thebibliography}{99}

\bibitem{FC}
M.~Firsching, P.~Lezeau, S.~Mercuri, M.~Z. Horv\'ath, Y.~Dillies, C.~S\"onne, E.~Wieser, F.~Zhang, T.~Hubert, B.~Ag\"uera y Arcas, and P.~Kohli, \emph{Formal Conjectures: An Open and Evolving Benchmark for Verified Discovery in Mathematics}, arXiv:2605.13171 [cs.AI], 2026. Repository: \url{https://github.com/google-deepmind/formal-conjectures}.

\bibitem{Mathlib}
The mathlib Community, \emph{The Lean mathematical library}, in: Proceedings of the 9th ACM SIGPLAN International Conference on Certified Programs and Proofs (CPP 2020), ACM, 2020, 367--381.

\bibitem{Lean4}
L.~de Moura and S.~Ullrich, \emph{The Lean 4 theorem prover and programming language}, in: Automated Deduction -- CADE 28, Lecture Notes in Comput. Sci. \textbf{12699}, Springer, 2021, 625--635.

\bibitem{A004086}
The OEIS Foundation Inc., \url{https://oeis.org/A004086} (read $n$ backwards), entry by N.~J.~A. Sloane.

\bibitem{A097153}
The OEIS Foundation Inc., \url{https://oeis.org/A097153} (reversal of $n$ minus largest digit of $n$), entry by J.~Earls (August~1, 2004).

\bibitem{A097154}
The OEIS Foundation Inc., \url{https://oeis.org/A097154}, entry by J.~Earls (August~1, 2004); revision \#4, December~15, 2017.

\bibitem{A097155}
The OEIS Foundation Inc., \url{https://oeis.org/A097155}, entry by J.~Earls (August~1, 2004), corrected by H.~P. Dale (2011); revision \#11, August~29, 2025.

\bibitem{A097156}
The OEIS Foundation Inc., \url{https://oeis.org/A097156}, entry by J.~Earls (August~1, 2004), more terms by R.~G. Wilson~v (2004).

\bibitem{A097157}
The OEIS Foundation Inc., \url{https://oeis.org/A097157}, entry by J.~Earls (August~1, 2004); revision \#3, December~15, 2017.

\bibitem{A097158}
The OEIS Foundation Inc., \url{https://oeis.org/A097158}, entry by J.~Earls (August~1, 2004), terms added by K.~Brockhaus (2004).

\end{thebibliography}
}

\end{document}
